\scanpage{51}
% archive sheet number: 50
\uline{Lemme}\quad Soit $\unsure{\mathcal{A}|A}$ un anneau d'endomorphismes
du \unsure{groupe|grpe|gpe} \deleted{\unsure{a}}abélien $H^{*}(X)$, stable par composantes homogènes
(degré \uuline{total}), par transposition. \deleted{\unsure{et tel}} \inserted{Supposons $\unsure{\xi|\zeta}\in\mathcal{A}$}, et que pour
tout $0\leq i<n$, il existe $u_i\in\mathcal{A}$, \inserted{de degré $\deleted{\unsure{2n}}-2(n-i)$,} tel que
$u_i\xi^{n-i}:H^i(X)\longrightarrow H^{2n-i}(X)\longrightarrow H^{\deleted{\unsure{2}}i}(X)$ soit l'identité.
Alors $\mathcal{A}$ contient $L$\margin{(\unsure{et réciproquement|et réciproquemt|et vice versa})}, \deleted{\unsure{les et}} et les projections
% slanted
sur les facteurs \unsure{élémentaires} $\xi^j P^i$ ($i\leq n$, et
\deleted{\unsure{et} $i+2j\leq n$ \unsure{ou} $i+2j\leq 2n-i$ \unsure{i.e.}} $j\leq n-i$)

\uline{Corollaire}\quad Conditions équivalentes :

(i) $L$ algébrique

(ii) pour \unsure{tt|et} $1\leq i\unsure{<|\leq}n$, il existe une corr. \uline{alg}
induisant $H^{2n-i}(X)\unsure{-|\to}H^i(X)$ inverse de $\unsure{J|f|\zeta}^i$

\deleted{\unsure{(iii)}} S'il en est ainsi, alors les \deleted{\unsure{$\xi$}} projections
sur les facteurs \unsure{élémentaires} de $H^{*}(X)$
sont algébriques, et \unsure{les} projections $\deleted{\unsure{de}}H^{*}(X)\unsure{\to|-}H^i(X)$
aussi.

\uline{\unsure{Dém|Dém.} du lemme.}\quad Comme $\mathcal{A}$ stable par
transposition, il suffit de regarder les $\unsure{\xi^jP^i|\zeta^jP^i}$,
pour $i+2j\leq n$. Soit $k=i+2j$, on procède
par récurrence sur $k$, trivial si $k<0$.

